\documentclass[11pt, a4paper]{article}

% bfw-style.tex — gemeinsame Style-Definitionen für alle Handouts/Aufgabenhefte
% Voraussetzung: Kompilation mit xelatex (wegen fontspec)

\usepackage[ngerman]{babel}
\usepackage{geometry}
\geometry{a4paper, top=2.2cm, bottom=2.2cm, left=2.3cm, right=2.3cm}

\usepackage{fontspec}
% Fallback-Kette: Source Sans Pro -> Calibri -> Segoe UI -> Latin Modern Sans
% (Latin Modern ist immer mit MiKTeX vorhanden)
\IfFontExistsTF{Source Sans Pro}{%
  \setmainfont{Source Sans Pro}[Ligatures=TeX]%
}{\IfFontExistsTF{Calibri}{%
  \setmainfont{Calibri}[Ligatures=TeX]%
}{\IfFontExistsTF{Segoe UI}{%
  \setmainfont{Segoe UI}[Ligatures=TeX]%
}{%
  \setmainfont{Latin Modern Sans}[Ligatures=TeX]%
}}}
\IfFontExistsTF{Source Code Pro}{%
  \setmonofont{Source Code Pro}[Scale=0.92]%
}{\IfFontExistsTF{Consolas}{%
  \setmonofont{Consolas}[Scale=0.92]%
}{%
  \setmonofont{Latin Modern Mono}[Scale=0.92]%
}}

\usepackage{xcolor}
% Farbpalette: hell, freundlich, modern
\definecolor{bfwPrimary}{HTML}{0EA5A4}    % Petrol/Türkis - Primär
\definecolor{bfwPrimaryDark}{HTML}{0F766E} % dunkler Petrol für Headings
\definecolor{bfwAccent}{HTML}{F59E0B}     % Amber - Akzent
\definecolor{bfwSuccess}{HTML}{10B981}    % Grün - Erfolg / Lösung
\definecolor{bfwWarn}{HTML}{F97316}       % Orange - Achtung
\definecolor{bfwInk}{HTML}{0F172A}        % fast Schwarz
\definecolor{bfwInkSoft}{HTML}{475569}    % gedämpftes Grau
\definecolor{bfwBgSoft}{HTML}{F8FAFC}     % off-white
\definecolor{bfwBgTeal}{HTML}{ECFEFF}     % helles Türkis
\definecolor{bfwBgAmber}{HTML}{FEF3C7}    % helles Gelb
\definecolor{bfwBgRose}{HTML}{FFE4E6}     % helles Rosé für Eselsbrücken

\usepackage[colorlinks=true, linkcolor=bfwPrimaryDark, urlcolor=bfwPrimaryDark]{hyperref}
\usepackage{enumitem}
\setlist[itemize]{leftmargin=*, itemsep=2pt, topsep=2pt}
\setlist[enumerate]{leftmargin=*, itemsep=2pt, topsep=2pt}

\usepackage[most]{tcolorbox}
\usepackage{booktabs}
\usepackage{tikz}
\usetikzlibrary{decorations.pathmorphing, positioning, shapes.geometric, arrows.meta}

% Merkkästchen — wichtige Definition / Kernpunkt
\newtcolorbox{merksatz}[1][]{%
  enhanced, breakable,
  colback=bfwBgTeal,
  colframe=bfwPrimary,
  boxrule=0.6pt,
  arc=4pt,
  left=10pt, right=10pt, top=8pt, bottom=8pt,
  fonttitle=\bfseries\color{bfwPrimaryDark},
  title={\faLightbulb\ Merksatz},
  #1
}

% Eselsbrücke — Mnemoniks
\newtcolorbox{eselsbruecke}[1][]{%
  enhanced, breakable,
  colback=bfwBgRose,
  colframe=bfwAccent,
  boxrule=0.6pt,
  arc=4pt,
  left=10pt, right=10pt, top=8pt, bottom=8pt,
  fonttitle=\bfseries\color{bfwWarn},
  title={Eselsbrücke},
  #1
}

% Beispiel-Box
\newtcolorbox{beispiel}[1][]{%
  enhanced, breakable,
  colback=bfwBgSoft,
  colframe=bfwInkSoft,
  boxrule=0.4pt,
  arc=3pt,
  left=10pt, right=10pt, top=8pt, bottom=8pt,
  fonttitle=\bfseries\color{bfwInk},
  title={Beispiel},
  #1
}

% Lösungs-Box (für Aufgabenhefte)
\newtcolorbox{loesung}[1][]{%
  enhanced, breakable,
  colback=bfwBgAmber!40,
  colframe=bfwSuccess,
  boxrule=0.6pt,
  arc=3pt,
  left=10pt, right=10pt, top=8pt, bottom=8pt,
  fonttitle=\bfseries\color{bfwSuccess!70!black},
  title={Lösung},
  #1
}

% Glossar-Eintrag
\newcommand{\glossareintrag}[2]{%
  \par\noindent\textbf{\color{bfwPrimaryDark}#1} \enspace #2\par\smallskip
}

% Section-Stil — etwas farbiger als Standard
\usepackage{titlesec}
\titleformat{\section}
  {\Large\bfseries\color{bfwPrimaryDark}}
  {\thesection}{0.6em}{}
\titleformat{\subsection}
  {\large\bfseries\color{bfwPrimary}}
  {\thesubsection}{0.5em}{}
\titleformat{\subsubsection}
  {\normalsize\bfseries\color{bfwInk}}
  {\thesubsubsection}{0.4em}{}

% TikZ-Stil "skizzenhaft" — etwas weicher als Rough.js, aber stabil
\tikzset{
  roughbox/.style = {
    draw=bfwPrimaryDark, thick, fill=bfwBgTeal,
    rounded corners=4pt, inner sep=6pt
  },
  roughaccent/.style = {
    draw=bfwAccent, thick, fill=bfwBgAmber,
    rounded corners=4pt, inner sep=6pt
  },
  roughline/.style = {
    draw=bfwInkSoft, thick, -{Stealth[length=2.5mm]}
  }
}

% Bullet-Symbole (bewusst kein FontAwesome — vermeidet Font-Installations-Probleme)
\newcommand{\faLightbulb}{$\bullet$}
\newcommand{\faBookmark}{$\bullet$}

% Header/Footer
\usepackage{fancyhdr}
\pagestyle{fancy}
\fancyhf{}
\renewcommand{\headrulewidth}{0pt}
\fancyfoot[L]{\small\color{bfwInkSoft}\thepage}
\fancyfoot[R]{\small\color{bfwInkSoft} BFW M-064 Beschaffung im Büromanagement}


\title{\vspace*{-1.5cm}\Huge\bfseries\color{bfwPrimaryDark}Aufgabenheft Tag~3\\[0.4em]
  \large\color{bfwInkSoft}\textnormal{Arbeitsumgebung \& Umweltfaktoren --- Übungen im IHK-Klausurformat}}
\author{\textsf{Büroprozesse gestalten und Arbeitsvorgänge organisieren \textperiodcentered{} M-067}\\
\textsf{\small\copyright\ Can Siebert\,\textperiodcentered\,2026}}
\date{Selbstlernmaterial \textperiodcentered{} 01.07.2026}

\begin{document}
\maketitle
\thispagestyle{empty}

\begin{merksatz}[title={\bfseries So nutzen Sie dieses Aufgabenheft}]
Bearbeiten Sie alle Aufgaben zunächst \emph{ohne} in die Lösungen zu schauen. Schreiben
Sie Ihre Antworten in vollständigen Sätzen --- die IHK-Prüfung verlangt formulierte
Begründungen, nicht bloße Stichworte. Beachten Sie die \emph{Operatoren} (nennen,
erläutern, beurteilen, begründen) --- sie geben den geforderten Antwort-Tiefegrad vor.
Erst danach öffnen Sie den Lösungsteil ab Seite~\pageref{loesungen} und vergleichen.
\end{merksatz}

\tableofcontents
\vspace{1em}
\hrule

\section{Aufgaben}

\subsection*{Aufgabe~1 --- Raumklima im Büro (12 Punkte)}

\noindent\textbf{\color{bfwAccent}YouTube-Vorbereitung:}\enspace \href{https://www.youtube.com/watch?v=IFWeO1r2Xj0}{Raumklima am Arbeitsplatz} (\url{https://www.youtube.com/watch?v=IFWeO1r2Xj0})

\medskip
Die \emph{Vertrieb Nord GmbH} klagt im Sommer über zu warme, im Winter über zu trockene
Büros. Die Geschäftsführung bittet Sie um eine fachliche Einschätzung.

\begin{enumerate}[label=(\alph*)]
  \item \textbf{Nennen} Sie die drei Bestandteile des Raumklimas. \hfill (3~P.)
  \item \textbf{Erläutern} Sie die Temperaturvorgaben der ASR A3.5 (20 / 26 / 30\,\textdegree C). \hfill (4~P.)
  \item \textbf{Nennen} Sie die optimale relative Luftfeuchtigkeit und die Folgen zu hoher bzw.\ zu niedriger Werte. \hfill (3~P.)
  \item \textbf{Beurteilen} Sie die Aussage: \emph{„Zugluft ist immer angenehm, weil sie kühlt.”} \hfill (2~P.)
\end{enumerate}

\subsection*{Aufgabe~2 --- Lüften und Klimatisierung (10 Punkte)}

\noindent\textbf{\color{bfwAccent}YouTube-Vorbereitung:}\enspace \href{https://www.youtube.com/watch?v=2m_39bPykSI}{Richtig lüften im Büro} (\url{https://www.youtube.com/watch?v=2m_39bPykSI})

\medskip
\begin{enumerate}[label=(\alph*)]
  \item \textbf{Erläutern} Sie, wann man von einer \emph{Klimatisierung} sprechen darf. \hfill (3~P.)
  \item \textbf{Nennen} Sie zwei Nachteile von Klimaanlagen. \hfill (2~P.)
  \item \textbf{Beschreiben} Sie, wie das Raumklima in nicht klimatisierten Räumen beeinflusst werden sollte. \hfill (3~P.)
  \item \textbf{Beurteilen} Sie, warum stündliches Stoßlüften wirksamer ist als ein dauerhaft gekipptes Fenster. \hfill (2~P.)
\end{enumerate}

\subsection*{Aufgabe~3 --- Beleuchtung am Bildschirmarbeitsplatz (12 Punkte)}

\noindent\textbf{\color{bfwAccent}YouTube-Vorbereitung:}\enspace \href{https://www.youtube.com/watch?v=Roz7ewsjEPA}{Beleuchtung am Bildschirm} (\url{https://www.youtube.com/watch?v=Roz7ewsjEPA})

\medskip
\begin{enumerate}[label=(\alph*)]
  \item \textbf{Nennen} Sie die geforderte Mindestbeleuchtungsstärke (in Lux) für Bildschirmarbeitsplätze. \hfill (2~P.)
  \item \textbf{Erläutern} Sie, warum eine gleichmäßig helle, flimmerfreie Beleuchtung wichtig ist. \hfill (3~P.)
  \item \textbf{Erklären} Sie den Unterschied zwischen \emph{Reflexblendung} und \emph{Direktblendung}. \hfill (3~P.)
  \item \textbf{Beurteilen} Sie, wie ein Sachbearbeiter mit Spiegelungen auf dem Bildschirm seinen Arbeitsplatz umgestalten sollte. \hfill (4~P.)
\end{enumerate}

\subsection*{Aufgabe~4 --- Lärm und Akustik (15 Punkte)}

\noindent\textbf{\color{bfwAccent}YouTube-Vorbereitung:}\enspace \href{https://www.youtube.com/watch?v=Sbj9Zk_Pg7M}{Lärm \& Schallschutz im Büro} (\url{https://www.youtube.com/watch?v=Sbj9Zk_Pg7M})

\medskip
In einem Großraumbüro mit 40 Plätzen klagen die Beschäftigten über Lärm durch Gespräche,
Telefonate und Druckergeräusche.

\begin{enumerate}[label=(\alph*)]
  \item \textbf{Erläutern} Sie, in welcher Einheit Lärm gemessen wird und was dieser Wert angibt. \hfill (3~P.)
  \item \textbf{Nennen} Sie die BAuA-Bewertung des Schalldruckpegels am Bildschirmarbeitsplatz (optimal / sehr gut / gut) sowie den Grenzwert für Gehörschutz. \hfill (4~P.)
  \item \textbf{Nennen} Sie vier Maßnahmen zur Lärmminderung im Büro. \hfill (4~P.)
  \item \textbf{Beurteilen} Sie, mit welchen modernen Akustiklösungen sich dieses Großraumbüro verbessern lässt. \hfill (4~P.)
\end{enumerate}

\subsection*{Aufgabe~5 --- Farbgestaltung und Sicherheitszeichen (12 Punkte)}

\noindent\textbf{\color{bfwAccent}YouTube-Vorbereitung:}\enspace \href{https://www.youtube.com/watch?v=aKkRtB-DZtE}{Farbpsychologie im Büro} (\url{https://www.youtube.com/watch?v=aKkRtB-DZtE})

\medskip
\begin{enumerate}[label=(\alph*)]
  \item \textbf{Erläutern} Sie die Wirkung warmer und kalter Farben sowie deren Einfluss auf die Raumwahrnehmung. \hfill (4~P.)
  \item \textbf{Nennen} Sie zwei positive Auswirkungen einer sinnvollen Farbverwendung in Büroräumen. \hfill (2~P.)
  \item \textbf{Ordnen} Sie den vier Sicherheitszeichen die richtige Farbe zu (Verbot, Warnung, Gebot, Rettung). \hfill (4~P.)
  \item \textbf{Beurteilen} Sie, warum Farbe und Beleuchtung gemeinsam betrachtet werden müssen. \hfill (2~P.)
\end{enumerate}

\subsection*{Aufgabe~6 --- Pflanzen, Strahlung und Gefahrstoffe (10 Punkte)}

\noindent\textbf{\color{bfwAccent}YouTube-Vorbereitung:}\enspace \href{https://www.youtube.com/watch?v=mU6sem1erMk}{Pflanzen im Büro} (\url{https://www.youtube.com/watch?v=mU6sem1erMk})

\medskip
\begin{enumerate}[label=(\alph*)]
  \item \textbf{Nennen} Sie drei positive Wirkungen von Pflanzen im Büro. \hfill (3~P.)
  \item \textbf{Nennen} Sie zwei Maßnahmen zur Reduktion elektromagnetischer Strahlung. \hfill (2~P.)
  \item \textbf{Erläutern} Sie, was Tonerstaub ist und welche Gesundheitsrisiken er bergen kann. \hfill (3~P.)
  \item \textbf{Beurteilen} Sie zwei sinnvolle Vorsichtsmaßnahmen im Umgang mit Tonerstaub. \hfill (2~P.)
\end{enumerate}

\vspace{1em}
\noindent\textbf{Punkte gesamt: 71}

\newpage

\section{Lösungen}\label{loesungen}

\subsection*{Lösung Aufgabe~1}
\begin{loesung}
\textbf{(a)} Die drei Bestandteile des Raumklimas: \textbf{Lufttemperatur}, \textbf{relative
Luftfeuchtigkeit} und \textbf{Luftgeschwindigkeit} (Luftbewegung).

\textbf{(b)} In Büroräumen soll eine Mindesttemperatur von \textbf{+20\,\textdegree C}
vorherrschen. Arbeitsräume \emph{sollen} \textbf{+26\,\textdegree C} nicht überschreiten ---
geschieht dies (z.\,B.\ durch Sonneneinstrahlung), sind Maßnahmen wie Jalousien, Gleitzeit,
Getränke oder gelockerte Bekleidungsregeln zu treffen. Werden \textbf{+30\,\textdegree C}
überschritten, \emph{müssen} zusätzliche Maßnahmen erfolgen.

\textbf{(c)} Optimal sind \textbf{40--60\,\%} relative Luftfeuchtigkeit. Zu hohe Werte
(v.\,a.\ im Winter) führen zu Schimmelgefahr; zu niedrige Werte führen zu brennenden Augen,
trockenen Lippen und Schleimhäuten und erhöhen die Infektionsgefahr.

\textbf{(d)} Die Aussage ist falsch. Nur eine geringe Luftbewegung von 0,1--0,15 m/s wird
als angenehm empfunden. Echte \emph{Zugluft} dagegen kann zu Nacken- und Rückenschmerzen
führen --- sie ist also gesundheitlich problematisch, nicht angenehm.
\end{loesung}

\subsection*{Lösung Aufgabe~2}
\begin{loesung}
\textbf{(a)} Von Klimatisierung darf nur gesprochen werden, wenn \emph{alle drei
Faktoren} (Lufttemperatur, Luftfeuchtigkeit, Luftbewegung) \textbf{automatisch geregelt}
werden.

\textbf{(b)} Nachteile von Klimaanlagen: Sie verbrauchen sehr viel Strom und können bei
schlechter Wartung Krankheitserreger übertragen.

\textbf{(c)} In nicht klimatisierten Räumen sollten die Fenster \emph{stündlich zum
Stoßlüften} geöffnet werden; außerdem bietet sich die Pausenzeit zum Lüften an. Das
Raumklima lässt sich zusätzlich durch Heizen, Kühlen oder Entlüften beeinflussen.

\textbf{(d)} Beim Stoßlüften wird in kurzer Zeit die gesamte verbrauchte Luft gegen
frische ausgetauscht, ohne die Wände auskühlen zu lassen. Ein dauerhaft gekipptes Fenster
tauscht dagegen kaum Luft aus, kühlt die Wände aus (Schimmelgefahr) und verschwendet
Heizenergie.
\end{loesung}

\subsection*{Lösung Aufgabe~3}
\begin{loesung}
\textbf{(a)} Gefordert sind \textbf{mindestens 500 Lux} für Bildschirmarbeitsplätze,
Arbeits- und Besprechungsräume, ergänzt durch persönliche Schreibtischleuchten.

\textbf{(b)} Ist der Raum gleichmäßig hell, müssen sich die Augen nicht ständig neu
einstellen --- das vermeidet Anstrengung, Ermüdung und Kopfschmerzen. Flimmerfreiheit
verhindert zusätzliche Belastung und Unruhe der Augen.

\textbf{(c)} \emph{Reflexblendung} entsteht durch Spiegelungen, z.\,B.\ auf dem Bildschirm;
\emph{Direktblendung} entsteht durch Lichtquellen im Blickfeld, z.\,B.\ Lampen oder das
Fenster. Beide erschweren das Sehen und lassen die Augen schneller ermüden.

\textbf{(d)} Der Bildschirm sollte \emph{parallel zum Fenster} aufgestellt werden, damit
sich das Fenster nicht darin spiegelt. Ergänzend helfen Sonnenschutz/Verschattung
(Jalousien) gegen direkte Einstrahlung, blendfreie Leuchten mit mindestens 500 Lux sowie
eine zusätzliche Schreibtischleuchte für Arbeiten auf Papier.
\end{loesung}

\subsection*{Lösung Aufgabe~4}
\begin{loesung}
\textbf{(a)} Lärm wird in \textbf{Dezibel --- dB(A)} --- gemessen; der Wert gibt den
\emph{Schalldruckpegel} an, also wie laut ein Geräusch ist.

\textbf{(b)} Laut BAuA gilt der Schalldruckpegel am Bildschirmarbeitsplatz bis
\textbf{30 dB(A)} als optimal, bis \textbf{40 dB(A)} als sehr gut und bis
\textbf{45 dB(A)} als gut. Ab \textbf{85 dB(A)} muss den Beschäftigten geeigneter,
persönlicher Gehörschutz bereitgestellt und auch getragen werden.

\textbf{(c)} Maßnahmen (vier von): geräuscharme Geräte auswählen; laute Geräte in
separaten Räumen unterbringen; Schallschutzhauben anschaffen; Schalldämpfung durch
Teppiche, Vorhänge oder schallabsorbierende Deckenplatten.

\textbf{(d)} Moderne Akustiklösungen: \emph{Akustikkabinen} (Telefon-/Fokuszellen) für
Telefonate und Konzentration, \emph{Schallabsorber-Stellwände} zwischen den
Arbeitsplätzen, \emph{Akustikdecken} und definierte Rückzugszonen. So sinkt der
Geräuschpegel, ohne die Vorteile der offenen Fläche aufzugeben.
\end{loesung}

\subsection*{Lösung Aufgabe~5}
\begin{loesung}
\textbf{(a)} \emph{Warme Farben} (Rot, Orange, Gelb) wirken aktivierend und stimulierend
und lassen einen Raum kleiner wirken. \emph{Kalte Farben} (Blau, Grün, Violett) wirken
beruhigend und distanzierend und lassen einen Raum größer wirken. Je kräftiger die Farbe,
desto stärker die Wirkung.

\textbf{(b)} Positive Auswirkungen (zwei von): Steigerung der Motivation; positive
Beeinflussung der Wahrnehmung ungünstiger Arbeitsbedingungen; Senkung der Energiekosten
(Beeinflussung des Temperaturempfindens); Unterstützung der Sicherheitskennzeichnung;
Förderung der Erholung.

\textbf{(c)} Zuordnung: \textbf{Verbotszeichen $=$ Rot}, \textbf{Warnzeichen $=$ Gelb},
\textbf{Gebotszeichen $=$ Blau}, \textbf{Rettungszeichen $=$ Grün} (Brandschutzzeichen
ebenfalls Rot).

\textbf{(d)} Farbe und Beleuchtung beeinflussen sich gegenseitig in ihrer Wirkung: Eine
Farbe wirkt unter warmem Licht anders als unter tageslichtweißem Licht. Wer Farben plant,
muss daher die Beleuchtung mitdenken, sonst entsteht ein unerwünschter Gesamteindruck.
\end{loesung}

\subsection*{Lösung Aufgabe~6}
\begin{loesung}
\textbf{(a)} Wirkungen von Pflanzen (drei von): sorgen für Entspannung und Wohlbefinden;
reinigen die Luft von Schadstoffen (z.\,B.\ Formaldehyd, Benzol, Kohlendioxid) und
Chemikalien; produzieren Sauerstoff; erhöhen die Luftfeuchtigkeit; schlucken Schall und
Staub (senken den Lärmpegel, beugen Atemwegserkrankungen vor).

\textbf{(b)} Maßnahmen gegen Strahlung (zwei von): strahlungsarme Geräte verwenden
(Flachbildschirme statt Röhrenbildschirme); im Büro auf tragbare Telefone verzichten;
nicht gebrauchte Geräte abschalten; W-LAN-Access-Points mindestens 3\,m vom Arbeitsplatz
entfernt anbringen.

\textbf{(c)} Tonerstaub gelangt aus Laserdruckern, Laserfax- und -kopiergeräten in die
Luft (z.\,B.\ beim Kartuschenwechsel oder Beseitigen eines Papierstaus). Mögliche Folgen:
schwere Atemwegserkrankungen, Hautprobleme, Reizung der Schleimhäute, Reizhusten,
Dauerschnupfen, Augenbrennen; manche Toner enthalten krebserregende Stoffe.

\textbf{(d)} Sinnvolle Vorsichtsmaßnahmen (zwei von): Geräte mit DGUV-Testzeichen und
schadstoffarme Toner anschaffen; Geräte in separaten, gut belüfteten Räumen aufstellen;
verschütteten Toner feucht aufwischen; nach Papierstaubeseitigung Hände waschen;
regelmäßige Wartung durch Fachpersonal; Kartuschen nicht gewaltsam öffnen.
\end{loesung}

\end{document}
